
\begin{frame}{Effet de l'analyse}

Quelques requêtes à faire sur la base des films au préalable:

\begin{itemize}
  \item \texttt{Matrix}, \texttt{MATRIX} et \texttt{matrix}
  \item \texttt{title:Matrix} et \texttt{text:Matrix}
  \item \texttt{title:reloaded} et \texttt{text:reloaded}
\end{itemize}

\vfill

\begin{block}{Que se passe-t-il?}
Si vous avez bien suivi les explications sur le schéma vous devriez
comprendre ce qui se passe. Réflechissons ensemble...
\end{block}
\end{frame}

\begin{frame}{Rôle de l'analyse}

L'analyse des textes permet d'effectuer une forme de 
normalisation / unification pour être moins dépendant de la
forme du texte.

\begin{itemize}
  \item un document parle de loup même si on y trouve les formes
     "loups", "Loup", "louve", etc.
  \item un document parle de travail quellle que soit la forme
     du verbe "travailler" ou de ses variantes.
  \item Jusqu'où va-t-on? Traductions (loup = wolf = lupus) ? Synonymes
      (loup  = prédateur) ? Recherche en cours.
\end{itemize}

\vfill

Pour moins dépendre de la forme on applique une \alert{transformation}.
\vfill

\begin{block}{Très important}
La même transformation doit être appliquée aux documents \alert{et} à la
requête (ex. de la requête "MATRIX")
\end{block}
\end{frame}



\begin{frame}{Impact de l'analyse}

Bien comprendre.

\vfill

\alert{Plus on normalise, plus on diminue la précision}.

Car des mots distincts sont unifiés (cote, côte, côté, etc.)
\vfill


\alert{Plus on normalise, plus on améliore le rappel}.

\vfill

Car on met en correspondance les variantes d'un même mot, d'une même
signification (conjuquaisons d'un verbe).

\end{frame}

\begin{frame}{Les phases de l'analyse}

Avant d'insérer dans l'index, plusieurs phases.

\vfill

Important: identification de quelques méta-données (la langue), prise en compte du contexte
(quels documents pour quelle application). Puis, de manière  générale\,:

\begin{itemize}
\item \alert{Tokenization}\,: découpage du texte en ``mots''
\item \alert{Normalisation}\,: majuscules\,? acronymes\,? apostrophes\,? accents\,?

     Exemple\,: \textit{Windows} et \textit{window}, U.S.A vs USA, \emph{l'auditeur}  vs \emph{les auditeurs}.
     
     \item \alert{Stemming} (``racinisation''), \alert{lemmatization}
     
     Prendre la racine des mots pour éviter le biais des variations (auditer, auditeur, audition, etc.)
     
     \item  \alert{Stop words}, quels mots garder\,?
     
     Mots très courants peu informatifs (le, un à, de).
\end{itemize}

C'est de l'art et du réglage... Dans ce qui suit\,: introduction / sensibilisation aux problèmes.

\end{frame}


\section{Tokenisation et normalisation}

\begin{frame}{Identification de la Langue}

Comment trouver la langue d'un document\,?

\begin{itemize}

\item Méta information (dans l'entête HTTP p.e.): pas fiable du tout.

\item \alert{Par le jeu de caractères}, pas assez courant!

\figSlide{characters}

\pause

\structure{Respectivement}: Coréen, Japonais, Maldives, Malte, Islandais.

\item Par extension: \alert{séquences de caractères fréquents}, ($n$-grams)

\item Par techniques d'apprentissage (\alert{classifiers})

\end{itemize}

Des librairies font ça très bien (e.g., Tika, \url{http://tika.apache.org})
\end{frame}

\begin{frame}{Tokenisation}

\begin{beamerboxesrounded}{Principe}
Séparation du texte en \alert{tokens} (``mots'')
\end{beamerboxesrounded}

\smallskip


Pas du tout aussi facile qu'on le dirait\,!

\begin{itemize}

\item Dans certaines langues (Chinois, Japonais), les mots  \alert{ne sont pas}
séparés par des espaces.

\item Certaines langues s'écrivent de droite à gauche, de haut en bas.

\item Que faire (et de manière \alert{cohérente}) des acronymes, élisions, nombres, unités, URL, email, etc.

\item \alert{Mots composés}\,: les séparer en \emph{tokens} ou les regrouper en un seul\,?

\begin{enumerate}
   \item Anglais\,: \emph{hostname}, \emph{host-name} et \emph{host
name}, ...
  \item Français\,: Le Mans, aujourd'hui, pomme de terre, ...
  \item Allemand\,: \textit{Levensversicherungsgesellschaftsangestellter} (employé d’une
société d’assurance vie)
\end{enumerate}

 Que faire si l'utilisateur
cherche \emph{hostname} et qu'on a normalisé en \emph{host-name}\,?

\end{itemize}


Majuscules, ponctuation\,? Une solution simple est de normaliser (minuscules, pas de ponctuation).

\end{frame}

\begin{frame}{Exemple pour notre petit jeu de données}

On met en minuscules, on retire la ponctuation.

\vfill

\begin{tabular}{>{\color{structure.fg}}ll}
     $d_1$&le loup est dans la bergerie\\
     $d_2$&le loup et les trois petits cochons\\
     $d_3$&les moutons sont dans la bergerie\\
     $d_4$&spider cochon spider cochon il peut marcher au plafond\\
     $d_5$&un loup a mangé un mouton les autres loups sont restés dans la bergerie\\
     $d_6$&il y a trois moutons dans le pré et un mouton dans la gueule du loup\\
     $d_7$&le cochon est à 12 euros le kilo le mouton à  10 euros le kilo\\
     $d_8$&les trois petits loups et le grand méchant cochon\\
\end{tabular}

\vfill

On considère que l'espace est le séparateur de tokens.

\end{frame}


\section{Racinisation}

\begin{frame}{Stemming (racine), lemmatization}
\begin{beamerboxesrounded}{Principe}
\alert{Confondre} toutes les formes d'un même mot, ou de mots apparentés, en
une seule \alert{racine}.
\end{beamerboxesrounded}

\smallskip

\begin{description}

\item[Stemming Morphologique.] Retire les pluriels, marque de genre, conjugaisons, modes, etc.

\begin{itemize}

\item Très dépendant de la langue\,: \emph{geese} pluriel de \emph{goose}, \emph{mice} de \emph{mouse}

\item Difficile à séparer d'une analyse linguistique (``Les poules du couvent couvent'', ``la petite
    brise la glace''\,:  où est le verbe\,?)
\end{itemize}

\item[Stemming lexical] Fondre les termes proches lexicalement\,: ``politique, politicien, police (?)'' ou 
     ``université, universel, univers (?)'' 


\item[Stemming phonétique.]  Correction fautes de frappes, fautes orthographes

\end{description}

%\begin{itemize}
  %\item lemmatization : travail morphologique plus poussé
%\end{itemize}

\end{frame}

\begin{frame}{Exemple de stemming}

On retire les pluriels, on met le verbe à l'infinitif.

\vfill

\begin{tabular}{>{\color{structure.fg}}ll}
     $d_1$&le loup etre dans la bergerie\\
     $d_2$&le loup et les trois petit cochon\\
     $d_3$&les moutons etre dans la bergerie\\
     $d_4$&spider cochon spider cochon il pouvoir marcher au plafond\\
     $d_5$&un loup avoir manger un mouton les autres loups etre rester dans la bergerie\\
     $d_6$&il y avoir trois mouton dans le pre et un mouton dans la gueule du loup\\
     $d_7$&le cochon etre a 12 euro le kilo le mouton a 10 euro le kilo\\
     $d_8$&les trois petit loup et le grand mechant cochon\\
\end{tabular}

\end{frame}


\section{Mots vides, synonymes, cas particuliers...}


\begin{frame}{Suppression des \textit{Stop Words}}

\begin{beamerboxesrounded}{Principe}
On retire les mots porteurs d'une information faible afin de limiter le stockage.
\end{beamerboxesrounded}

\begin{description}
\item[articles:] \emph{le}, \emph{le}, \emph{ce}, etc.

\item[verbes ``fonctionnels''] \emph{être}, \emph{avoir}, \emph{faire}, etc.

\item[conjunctions:] \emph{that}, \emph{and}, etc.

\item[etc.]
\end{description}

\begin{remark}
Maintenant moins utilisé car (i) espace de stockage peu coûteux et (ii) pose d'autres problèmes (``pomme de terre'',
``Let it be'', ``Stade de France'')
\end{remark}

\end{frame}



\begin{frame}{Autres problèmes, en vrac}


\alert{Majuscules / minuscules}

Lyonnaise des Eaux, Société Générale, etc.

\vfill
 \alert{Acronymes}
 
   CAT = \textit{cat} ou \textit{Caterpillar Inc.}\,?  M.A.A.F ou MAAF ou Mutuelle ... \,?

\vfill

\alert{Dates, chiffres}

   Monday 24, August, 1572 -- 24/08/1572 -- 24 août 1572
    10000 ou 10,000.00 ou 10,000.00
     
     
     \vfill
     
     \alert{Accents, ponctuation}
     
     résumé ou résume ou resume... 
     
 \begin{remark}
 Dans tous les cas, les même règles de transformation s'appliquent aux documents ET à la requête.
 \end{remark}
 \end{frame}


\begin{frame}{Exemple avec suppression des \textit{stop words}}

Voici une solution possible.

\vfill

\begin{tabular}{>{\color{structure.fg}}ll}
     $d_1$&loup etre bergerie\\
     $d_2$&loup  trois petit cochon\\
     $d_3$&mouton etre bergerie\\
     $d_4$&spider cochon spider cochon pouvoir marcher plafond\\
     $d_5$&loup avoir manger mouton  autres loups etre rester bergerie\\
     $d_6$&avoir trois mouton pre mouton gueule loup\\
     $d_7$&cochon etre 12 euro kilo mouton 10 euro kilo\\
     $d_8$&trois petit loup grand mechant cochon\\
\end{tabular}

\vfill

On a gardé les verbes fonctionnels (être, avoir).
\end{frame}

\section{Analyse avec Solr}

\begin{frame}[fragile]{Définir une chaîne d'analyse avec Solr}

Un analyseur est associé à un champ:

\begin{itemize}
  \item le tokenizer effectue le traitement lexical consistant à transformer le texte en un ensemble de tokens ;
   \item les filtres (filter) examinent les tokens un par un et décident de les conserver, de les remplacer par un ou plusieurs autres;
   \item l’analyseur est une chaîne de traitement (pipeline) constituée de tokeniseurs et de filtres.
\end{itemize}

\vfill
Exemple minimal:

\begin{minted}[frame=lines,                                                     
               baselinestretch=.9,                                              
               fontsize=\tiny,                                          
               framesep=2mm]{xml}
               
   <fieldType name="text" class="solr.TextField">
		<analyzer>
	   	 <tokenizer class="solr.StandardTokenizerFactory" />
	     <filter class="solr.LowerCaseFilterFactory" />
	</analyzer>
   </fieldType>
\end{minted}
  
\end{frame}



\begin{frame}[fragile]{tester les analyseurs}

Très pratique: l'interface \url{http://localhost:8983/solr/#/movies/analysis}.

\vfill

\begin{itemize}
  \item      dans le champ field value (index), on saisit un texte à analyser;
   \item dans le menu déroulant au dessous, on choisit un 
     des champs du schéma de l’index; c’est donc l’analyseur de ce champ 
     qui sera appliqué.
\end{itemize}

\vfill

Démonstration !



\end{frame}

